\documentclass[12pt]{article}
\input{header.tex}

\begin{document}

\section{Билинейные формы}

\begin{definition}
Пусть \(L\) — конечномерное линейное пространство над \(\mathbb{R}\). Отображение \(B(x, y)\), составляющее каждой упорядоченной паре элементов \(x, y\) из \(L\) вещественное число, называется \textbf{билинейной формой}, заданной на \(L\), если выполнены условия
\begin{enumerate}
    \item \(\forall \alpha, \beta \in \mathbb{R},\; \forall x, y, z \in L\) имеем \(B(\alpha x + \beta y, z) = \alpha B(x, z) + \beta B(y, z)\) --- линейность по первому аргументу;
    \item \(\forall \alpha, \beta \in \mathbb{R},\; \forall x, y, z \in L\) имеем \(B(x, \alpha y + \beta z) = \alpha B(x, y) + \beta B(x, z)\) --- линейность по второму аргументу.
\end{enumerate}
\end{definition}

\begin{definition}
Билинейная форма \(B(x, y)\) называется \textbf{симметрической}, если \(\forall x, y \in L\) имеем \(B(x, y) = B(y, x)\).
\end{definition}

\begin{theorem}
Пусть \(e = (e_1, \dots, e_n)\) — базис \(L\), \(B(x, y)\) — билинейная форма на \(L\). Тогда существует единственный набор вещественных чисел \(b_{ij}\;(1 \le i, j \le n)\) такой, что при любых \(x\) и \(y\) имеем
\begin{equation} \label{eq:bilinear_gen}
B(x, y) = \sum_{i,j=1}^{n} b_{ij} x_i y_j,
\end{equation}
где \(x_i, y_j\) — координаты в базисе \(e\) векторов \(x\) и \(y\) соответственно. При этом \(b_{ij} = B(e_i, e_j)\).
\end{theorem}

\begin{proof}
Разложим векторы \(x\) и \(y\) по базису \(e\):
\[
x = \sum_{i=1}^{n} x_i e_i, \quad y = \sum_{j=1}^{n} y_j e_j.
\]
По линейности \(B(x, y)\) получаем
\[
B(x, y) = \sum_{i,j=1}^{n} B(e_i, e_j) x_i y_j.
\]
Осталось ввести обозначение \(b_{ij} = B(e_i, e_j)\), и существование доказано. Проверим единственность. Пусть
\[
B(x, y) = \sum_{i,j=1}^{n} b_{ij} x_i y_j.
\]
Подставив \(x = e_i, y = e_j\), видим, что \(b_{ij} = B(e_i, e_j)\).
\end{proof}

\subsection{Квадратичные формы}
\begin{definition}
\textbf{Квадратичной формой} $Q(x)$ на пространстве $V$ называется функция, которая в некотором базисе имеет вид однородного многочлена второй степени от координат вектора:
\[
Q(x) = \sum_{i=1}^n \sum_{j=1}^n a_{ij} x_i x_j.
\]
\end{definition}

С каждой квадратичной формой $Q(x)$ можно связать единственную симметрическую билинейную форму $B(x,y)$, называемую \textbf{полярной}, такую что $Q(x) = B(x,x)$. Она находится по формуле \textit{поляризации}:
\[
\boxed{B(x, y) = \frac{1}{2}(Q(x+y) - Q(x) - Q(y))}
\]

\begin{proof}[Доказательство формулы поляризации]
По определению, $Q(x) = B(x,x)$. Вычислим $Q(x+y)$:
\[
Q(x+y) = B(x+y, x+y).
\]
Используя билинейность и симметричность $B$, раскроем:
\[
B(x+y, x+y) = B(x,x) + B(x,y) + B(y,x) + B(y,y) = B(x,x) + 2B(x,y) + B(y,y).
\]
Подставляя $Q(x)$ и $Q(y)$, получаем:
\[
Q(x+y) = Q(x) + 2B(x,y) + Q(y).
\]
Отсюда $2B(x,y) = Q(x+y) - Q(x) - Q(y)$, что и требовалось.
\end{proof}

Если $Q(x) = X^T A X$, то $B(x,y) = X^T A Y$, причем матрица $A$ может быть выбрана симметричной ($A^T = A$). Везде далее подразумевается, что матрица квадратичной формы симметрична.

\begin{remark}
При замене базиса $X = C X'$, матрица квадратичной формы преобразуется как $A' = C^T A C$, сохраняя симметричность. Ранг матрицы $A$ называется \textbf{рангом квадратичной формы} $\rg Q$ и является инвариантом относительно невырожденных линейных преобразований.
\end{remark}

\section{Приведение к каноническому и нормальному виду}
Основная задача: найти невырожденное преобразование $X = CY$, приводящее квадратичную форму к \textit{каноническому виду}:
\[
Q(x) = \lambda_1 y_1^2 + \lambda_2 y_2^2 + \dots + \lambda_n y_n^2.
\]

\begin{definition}
Канонический вид называется \textbf{нормальным}, если все ненулевые коэффициенты $\lambda_i$ равны $+1$ или $-1$. Всегда можно перейти от канонического к нормальному виду, выполнив замену $y_i = z_i / \sqrt{|\lambda_i|}$ (для $\lambda_i \neq 0$).
\end{definition}

\subsection{Метод Лагранжа (выделения полных квадратов)}
Универсальный метод приведения к каноническому виду.

\begin{stat}[Метод Лагранжа]
Рассмотрим два случая:
\begin{enumerate}
    \item \textbf{Существует $a_{ii} \neq 0$}. Без ограничения общности, пусть $a_{11} \neq 0$. Сгруппируем все слагаемые, содержащие $x_1$:
    \[
    Q(x) = a_{11}x_1^2 + 2\sum_{j=2}^n a_{1j}x_1x_j + \tilde{Q}(x_2,\dots,x_n).
    \]
    Выделим полный квадрат:
    \[
    Q(x) = \frac{1}{a_{11}}\left( a_{11}x_1 + \sum_{j=2}^n a_{1j}x_j \right)^2 - \frac{1}{a_{11}} \left( \sum_{j=2}^n a_{1j}x_j \right)^2 + \tilde{Q}(x_2,\dots,x_n).
    \]
    Введем новую переменную $y_1 = \sum_{j=1}^n a_{1j}x_j$, и оставим $y_i = x_i$ для $i>1$. Это преобразование невырождено, так как $a_{11} \neq 0$. В результате:
    \[
    Q(x) = \frac{1}{a_{11}} y_1^2 + Q_1(y_2,\dots,y_n),
    \]
    где $Q_1$ --- квадратичная форма от $n-1$ переменной.
    \item \textbf{Все $a_{ii}=0$, но $\exists a_{ij}\neq0$}. Пусть $a_{12}\neq0$. Выполним вспомогательную замену:
    \[
    \begin{cases}
    x_1 = z_1 + z_2, \\
    x_2 = z_1 - z_2, \\
    x_k = z_k, \quad k>2.
    \end{cases}
    \]
    Тогда $2a_{12}x_1x_2 = 2a_{12}(z_1^2 - z_2^2)$, и в новой форме появятся ненулевые коэффициенты при квадратах. Далее применяем пункт 1.
\end{enumerate}
Повторяя процесс, после конечного числа шагов получим канонический вид.
\end{stat}

\subsection{Метод Якоби}
Метод применим, если все угловые определители $\Delta_k = \det A_k$ матрицы $A$ отличны от нуля ($A_k$ --- угловая подматрица матрицы $A$ размера $k$, то есть подматрица, составленная из первых $k$ строк и столбцов; $k = 1, \dots, n$).

\begin{theorem}[Якоби]
Пусть $\Delta_k \neq 0$ для всех $k = 1, \dots, n-1$. Тогда существует единственное невырожденное преобразование с верхней унитреугольной матрицей $Q$ (то есть $Q$ верхнетреугольная с единицами на диагонали) такое, что
\[
Q^T A Q = \operatorname{diag}\!\left( \frac{\Delta_1}{\Delta_0},\; \frac{\Delta_2}{\Delta_1},\; \dots,\; \frac{\Delta_n}{\Delta_{n-1}} \right), \quad \Delta_0 = 1.
\]
Столбцы $q_k$ матрицы $Q$ строятся следующим образом: $q_{kk} = 1$, элементы ниже диагонали равны нулю, а элементы выше диагонали $(q_{1k}, \dots, q_{k-1,k})^T$ определяются из системы
\[
A_{k-1} \begin{pmatrix} q_{1k} \\ \vdots \\ q_{k-1,k} \end{pmatrix}
= -\begin{pmatrix} a_{1k} \\ \vdots \\ a_{k-1,k} \end{pmatrix},
\]
имеющей единственное решение при $\det A_{k-1} = \Delta_{k-1} \neq 0$.
\end{theorem}

\begin{proof}
Разложение $Q^T A Q = D$ получается из $LDU$-разложения матрицы $A$ (родственного $LU$-разложению): для симметричной матрицы с ненулевыми угловыми определителями существует единственное разложение $A = LDL^T$, где $L$ --- нижняя унитреугольная, $D$ --- диагональная с элементами $\Delta_k / \Delta_{k-1}$, откуда $Q = (L^T)^{-1}$ является верхней унитреугольной и даёт нужное приведение. Корректность приведённых формул для нахождения столбцов матрицы $Q$ доказывается индукцией по размеру матрицы.
\end{proof}

\subsection{Ортогональное преобразование}
Рассмотрим случай $V = \R^n$ со стандартным скалярным произведением. Так как $A = A^T$, то для нее существует ортонормированный базис из собственных векторов. Матрица перехода к этому базису $Q$ является ортогональной ($Q^{-1} = Q^T$). Тогда
\[
Q^T A Q = \operatorname{diag}(\lambda_1, \dots, \lambda_n),
\]
где $\lambda_i$ --- собственные числа матрицы $A$. Это преобразование также приводит форму к каноническому виду, но коэффициентами являются собственные числа.

\section{Закон инерции}
\begin{theorem}[Закон инерции]
Невырожденное линейное преобразование не меняет сигнатуру квадратичной формы. То есть, если форма $f(x) = x^T A x$ приведена к каноническому виду двумя разными способами:
\[
g(y) = \lambda_1 y_1^2 + \dots + \lambda_p y_p^2 - \lambda_{p+1} y_{p+1}^2 - \dots - \lambda_{p+k} y_{p+k}^2,
\]
\[
t(z) = \mu_1 z_1^2 + \dots + \mu_s z_s^2 - \mu_{s+1} z_{s+1}^2 - \dots - \mu_{s+l} z_{s+l}^2,
\]
где $\lambda_i, \mu_j > 0$, то $p = s$ и $k = l$.
\end{theorem}

\begin{proof}
От противного. Предположим $p < s$. Обозначим преобразование к каноническому виду $g$ как $Q_1$ ($x = Q_1 y$), преобразование к каноническому виду $t$ как $Q_2$ ($x = Q_2 z$).

Рассмотрим систему из $p + (n - s)$ однородных уравнений:
\[
y_1 = 0,\quad \dots,\quad y_p = 0, \qquad z_{s+1} = 0,\quad \dots,\quad z_n = 0,
\]
связывающую координаты $x$ через $Q_1$ и $Q_2$. Поскольку $p < s$, имеем $p + (n - s) < n$, то есть число уравнений строго меньше числа неизвестных. Следовательно, существует ненулевой вектор $x_0$, удовлетворяющий этой системе.

Вычислим $f(x_0)$ двумя способами. Через первую каноническую форму (при $y_1 = \dots = y_p = 0$):
\[
f(x_0) = g(Q_1^{-1} x_0) = -\lambda_{p+1} y_{p+1}^2 - \dots - \lambda_{p+k} y_{p+k}^2 \;\le\; 0.
\]
Через вторую каноническую форму (при $z_{s+1} = \dots = z_n = 0$):
\[
f(x_0) = t(Q_2^{-1} x_0) = \mu_1 z_1^2 + \dots + \mu_s z_s^2 \;\ge\; 0.
\]
Следовательно, $f(x_0) = 0$. Из второго выражения, поскольку все $\mu_j > 0$:
\[
z_1 = \dots = z_s = 0.
\]
Вместе с $z_{s+1} = \dots = z_n = 0$ получаем $z = 0$, откуда $Q_2^{-1} x_0 = 0$. Так как $Q_2$ невырождена, $x_0 = 0$ --- противоречие.
\end{proof}

\begin{definition}
Набор $(\sigma^+, \sigma^-, \sigma^0)$, где $\sigma^+$ --- число положительных, $\sigma^-$ --- число отрицательных коэффициентов в каноническом виде, а $\sigma^0 = n - \rank A$ --- число нулей, называется \textbf{сигнатурой} квадратичной формы. Число $\sigma = \sigma^+ - \sigma^-$ также называют сигнатурой.
\end{definition}

\section{Знакоопределенные квадратичные формы}

\begin{definition}
Квадратичная форма $Q(x)$ на $\R^n$ называется:
\begin{itemize}
    \item \textbf{Положительно определенной} ($Q > 0$), если $Q(x) > 0$ для всех $x \neq 0$.
    \item \textbf{Отрицательно определенной} ($Q < 0$), если $Q(x) < 0$ для всех $x \neq 0$.
    \item \textbf{Положительно полуопределенной} ($Q \ge 0$), если $Q(x) \ge 0$ для всех $x$.
    \item \textbf{Отрицательно полуопределенной} ($Q \le 0$), если $Q(x) \le 0$ для всех $x$.
    \item \textbf{Неопределенной}, если она принимает как положительные, так и отрицательные значения.
\end{itemize}
\end{definition}

\subsection{Критерии знакоопределенности}

\begin{theorem}[Критерий в терминах собственных чисел]
Для квадратичной формы $Q(x) = X^T A X$:
\begin{align*}
Q > 0 &\iff \text{все собственные числа } \lambda_i(A) > 0. \\
Q < 0 &\iff \text{все собственные числа } \lambda_i(A) < 0. \\
Q \ge 0 &\iff \text{все собственные числа } \lambda_i(A) \ge 0. \\
Q \le 0 &\iff \text{все собственные числа } \lambda_i(A) \le 0.
\end{align*}
\end{theorem}

\begin{theorem}[Критерий в терминах сигнатуры]
Пусть $\rg Q = r$.
\begin{align*}
Q > 0 &\iff \sigma^+ = n, \sigma^- = 0. \\
Q < 0 &\iff \sigma^+ = 0, \sigma^- = n. \\
Q \ge 0 &\iff \sigma^- = 0. \\
Q \le 0 &\iff \sigma^+ = 0. \\
Q \text{ неопределена} &\iff \sigma^+ > 0 \text{ и } \sigma^- > 0.
\end{align*}
\end{theorem}

\begin{proof}[Доказательство критерия через сигнатуру]
Приведем форму к нормальному виду (существует невырожденное преобразование):
\[
Q(x) = y_1^2 + \dots + y_p^2 - y_{p+1}^2 - \dots - y_r^2,
\]
где $p=\sigma^+$, $m=\sigma^-$, $r=p+m$, а нулевые члены опущены. Поскольку преобразование невырождено, множеству $x\neq0$ взаимно-однозначно соответствует множество $y\neq0$. Исследуем знак $Q$ в координатах $y$.

1. $Q>0$: Если $m>0$, взяв $y$ с $y_{p+1}=1$, получим $Q=-1<0$ — противоречие. Значит $m=0$. Если $r<n$ (т.е. есть нулевые координаты), взяв вектор с единственной ненулевой координатой $y_n=1$, получим $Q=0$ при $x\neq0$ — противоречие. Значит $r=n$, т.е. $\sigma^+=n$, $\sigma^-=0$.

2. $Q<0$: Аналогично, $p=0$, $m=n$.

3. $Q\ge0$: Если $m>0$, снова найдется вектор с $Q<0$. Следовательно, $m=0$. Наличие нулевых членов допускается.

4. $Q\le0$: Аналогично, $p=0$.

5. Неопределенность: $p>0$ и $m>0$ очевидно дают и положительные, и отрицательные значения.
\end{proof}

\subsection{Критерий Сильвестра}
Наиболее практически важный критерий для проверки положительной определенности.

\begin{theorem}[Критерий Сильвестра]
Квадратичная форма $Q(x) = X^T A X$ является положительно определенной тогда и только тогда, когда все главные (угловые) миноры ее матрицы положительны:
\[
\Delta_1 = a_{11} > 0,\ \Delta_2 = \begin{vmatrix} a_{11} & a_{12} \\ a_{12} & a_{22} \end{vmatrix} > 0,\ \dots,\ \Delta_n = \det A > 0.
\]
\end{theorem}

\begin{proof}
\textit{Необходимость.} Если $Q>0$, то ее ограничение на подпространство $\langle e_1,\dots,e_k\rangle$ также положительно определено. Матрица этого ограничения есть $A_k$, и ее определитель положителен (как произведение положительных собственных чисел).
\textit{Достаточность.} Так как все $\Delta_k \neq 0$, применим метод Якоби. Существует базис, в котором форма имеет вид $Q = \Delta_1 y_1^2 + (\Delta_2/\Delta_1) y_2^2 + \dots + (\Delta_n/\Delta_{n-1}) y_n^2$. Так как все коэффициенты положительны, форма положительно определена.
\end{proof}

\begin{theorem}[Критерий Сильвестра для отрицательной определенности]
Форма $Q(x)$ отрицательно определена тогда и только тогда, когда знаки угловых миноров чередуются, начиная с минуса:
\[
\Delta_1 < 0,\ \Delta_2 > 0,\ \Delta_3 < 0,\ \dots,\ (-1)^n \Delta_n > 0.
\]
\end{theorem}
\begin{proof}
Применить критерий Сильвестра к форме $-Q$, матрица которой есть $-A$.
\end{proof}

\begin{corollary}
Пусть все угловые миноры квадратичной формы $Q(x) = X^T A X$: $\Delta_k \neq 0$ для $k=1..n-1$. В таком случае $Q$ является положительно полуопределенной ($Q \ge 0$), тогда и только тогда, когда все главные (угловые) миноры ее матрицы положительны, а определитель матрицы неотрицателен:
\[
\Delta_1 = a_{11} > 0,\ \Delta_2 = \begin{vmatrix} a_{11} & a_{12} \\ a_{12} & a_{22} \end{vmatrix} > 0,\ \dots,\ \Delta_n = \det A \geq 0.
\]
\end{corollary}

\begin{theorem}[Критерий полуопределённости]
Квадратичная форма $Q(x) = X^T A X$ является положительно полуопределенной ($Q \ge 0$) тогда и только тогда, когда все ее главные миноры (все, а не только угловые!) неотрицательны: $\Delta \ge 0$.
\end{theorem}

\begin{proof}
(Только для необходимости) Рассмотрим любое подпространство, натянутое на $k$ базисных векторов, например, $\langle e_{i_1}, \dots, e_{i_k} \rangle$. Ограничение формы на это подпространство также является положительно полуопределенной квадратичной формой. Матрица этого ограничения является соответствующей главной подматрицей $A_{i_1 \dots i_k}$. Для любой положительно полуопределенной матрицы все ее главные миноры неотрицательны (это следует из того, что все собственные числа такой матрицы $\ge 0$, а определитель равен их произведению).
\end{proof}

\begin{example}[Контрпример]
Рассмотрим матрицу $A = \begin{pmatrix} 0 & 0 \\ 0 & -1 \end{pmatrix}$. Ее угловые миноры: $\Delta_1 = 0 \ge 0$, $\Delta_2 = 0 \ge 0$. Однако соответствующая форма $Q(x,y) = -y^2$ является \textbf{отрицательно полуопределенной}, но не положительно полуопределенной. Здесь главный минор, соответствующий элементу $a_{22} = -1$, отрицателен.
\end{example}

\section*{Приложение: Поверхности второго порядка}
Квадратичные формы широко используются для классификации кривых и поверхностей второго порядка. Уравнение поверхности второго порядка в $\R^3$ имеет вид:
\[
a_{11}x^2 + a_{22}y^2 + a_{33}z^2 + 2a_{12}xy + 2a_{13}xz + 2a_{23}yz + 2b_1x + 2b_2y + 2b_3z + c = 0.
\]
В матричной форме: $X^T A X + 2B^T X + c = 0$, где $A = A^T$.

\begin{enumerate}
    \item \textbf{Ортогональное преобразование} $X = QX'$ приводит квадратичную часть к каноническому виду $\lambda_1 x'^2 + \lambda_2 y'^2 + \lambda_3 z'^2$, одновременно преобразуя линейную часть: $B^T Q X'$.
    \item \textbf{Выделение полных квадратов} (параллельный перенос) позволяет избавиться от линейных членов для тех переменных, у которых $\lambda_i \neq 0$. В результате уравнение приводится к одному из канонических видов (эллипсоид, гиперболоиды, параболоиды, цилиндры и т.д.).
\end{enumerate}

\end{document}